\section{Reproducibility}
\label{app:repro}

This appendix gives the exact command, model strings, seed manifest, and software environment needed to rerun the experiments of Section~\ref{sec:agents}. Per-panel results are reported separately in Appendix~\ref{app:results}; this section is restricted to what a reviewer needs to regenerate the run inputs.

\subsection{Reproduction command}
\label{app:repro:command}

The canonical invocation, anchored on the GPT-5.5 panel, is:

\begin{verbatim}
LADDER_MAX_RETRIES=20 uv run python run_ladder.py \
    --levels 0,1,2,3,4,5 \
    --seeds-per-level 8 \
    --models gpt-5.5 \
    --alpha 0.05 \
    --max-steps-raw 20 \
    --max-steps-stats 40 \
    --preflight-seed 20260422 \
    --out-dir traces/ladder/full_gpt55
\end{verbatim}

Substituting the \texttt{--models} string from the table in Appendix~\ref{app:repro:models} and changing \texttt{--out-dir} to the matching \texttt{Trace directory} listed there reproduces the corresponding panel. The env-var prefix \texttt{LADDER\_MAX\_RETRIES=20} enables an auto-resume loop (\texttt{run\_ladder.py:1289--1304}) that retries on transient provider errors, restoring from the on-disk checkpoint after a 10-second backoff. The default value is $0$ (no retry), so the env-var is required to match the canonical runs.

\subsection{Model versions}
\label{app:repro:models}

The model string and run identifier for each panel are lifted verbatim from \texttt{run\_manifest.json} in the corresponding output directory. All five runs share \texttt{levels=[0,1,2,3,4,5]}, \texttt{seeds\_per\_level=8}, $\alpha=0.05$, $\texttt{max\_steps\_raw}=20$, and $\texttt{max\_steps\_stats}=40$.

\begin{center}
\scriptsize
\setlength{\tabcolsep}{4.0pt}
\begin{tabular}{@{}l l l l l@{}}
\toprule
Panel & Model string & Provider & \texttt{run\_id} (UTC) & Trace directory \\
\midrule
GPT-5.5            & \texttt{gpt-5.5}                                 & OpenAI       & \texttt{20260503T183924Z} & \texttt{traces/ladder/full\_gpt55} \\
GPT-5.4-mini       & \texttt{gpt-5.4-mini}                            & OpenAI       & \texttt{20260504T054452Z} & \texttt{traces/ladder/full\_gpt54mini} \\
Claude Sonnet 4.6  & \texttt{openrouter/anthropic/claude-sonnet-4-6}  & OpenRouter   & \texttt{20260504T132941Z} & \texttt{traces/ladder/full\_sonnet46\_or} \\
Claude Haiku 4.5   & \texttt{openrouter/anthropic/claude-haiku-4.5}   & OpenRouter   & \texttt{20260504T111120Z} & \texttt{traces/ladder/full\_haiku45\_or} \\
Gemini 3 Flash     & \texttt{google/gemini-3-flash-preview}           & LiteLLM      & \texttt{20260504T172409Z} & \texttt{traces/ladder/full\_gemini3flash} \\
\bottomrule
\end{tabular}
\end{center}

Model strings without a provider prefix are dispatched by \texttt{provider\_for\_model} to the OpenAI or Anthropic native APIs; prefixed strings (\texttt{anthropic/...}, \texttt{openrouter/...}, \texttt{google/...}) are routed through LiteLLM. The Sonnet and Haiku panels use the \texttt{\_or} suffix to mark the canonical OpenRouter re-runs that supersede earlier partial runs at the same \texttt{run\_id}; the aggregated metrics in \texttt{traces/aggregated/per\_model\_per\_level\_per\_method.csv} are sourced exactly from the five directories listed above.

\subsection{Seed manifest}
\label{app:repro:seeds}

The same accepted-seed map is reused across all five panels --- it appears verbatim in every \texttt{run\_manifest.json} --- so cross-model comparisons are paired by instance.

\begin{itemize}
\itemsep0pt
\item \textbf{L0:} \texttt{1276453582, 27998312, 931551983, 974369856, 1539978433, 868078822, 452641024, 963592850}
\item \textbf{L1:} \texttt{1257633785, 1100146725, 1910407594, 746528573, 1175759265, 1982942393, 1449442515, 1479720101}
\item \textbf{L2:} \texttt{823847298, 834246677, 32015887, 716908331, 1426457537, 684342303, 1207468583, 1612783012}
\item \textbf{L3:} \texttt{1329552115, 160719912, 362855691, 589330817, 1877150515, 370845052, 976218778, 1286035187}
\item \textbf{L4:} \texttt{518645330, 511033471, 1198956094, 154587426, 1368642322, 1163629188, 1332933138, 1231061499}
\item \textbf{L5:} \texttt{1526642568, 1170708238, 385840006, 32447648, 1017236502, 1717602974, 480184089, 2015902496}
\end{itemize}

The map is generated by \texttt{build\_seed\_map\_preflight} (\texttt{run\_ladder.py:189--217}) with \texttt{--preflight-seed 20260422}, accepting a candidate seed only if it passes the rejection filters of Appendix~\ref{app:instance:filters}, then frozen and replayed via the manifest. Pinning \texttt{--preflight-seed} together with the filter spec is sufficient to regenerate the map from scratch; copying the 48 integers above is provided only for direct replay.

\subsection{Software environment}
\label{app:repro:env}

\begin{itemize}
\itemsep0pt
\item Python $\ge 3.12$, declared in \texttt{pyproject.toml}.
\item Install via \texttt{uv sync}, which materialises the locked dependency set.
\item LLM dispatch goes through \texttt{litellm} $\ge 1.79$ (with \texttt{openai} $\ge 2.30$ for native OpenAI calls). PC inference uses \texttt{causal-learn} $\ge 0.1.4.5$. Orchestration relies on \texttt{numpy}, \texttt{pandas} $\ge 3.0.2$, \texttt{python-dotenv} $\ge 1.2.2$, and \texttt{tenacity} $\ge 9.1.4$.
\item Provider credentials \texttt{OPENAI\_API\_KEY}, \texttt{ANTHROPIC\_API\_KEY}, and \texttt{OPENROUTER\_API\_KEY} are read from \texttt{.env} at startup (\texttt{run\_ladder.py:1069--1095}); the runner refuses to launch if a key required by the requested model panel is missing.
\end{itemize}

\subsection{Repository}
\label{app:repro:repo}

The full source, manifests, and per-panel \texttt{events.jsonl} traces are available at \url{https://anonymous.4open.science/r/ACDB}.
